\documentclass[12pt]{article}
\usepackage{geometry}
\geometry{margin=1in}

% ----------------------------------------------------------
%  Linux Penguin Theme — Based on Recommended Colors
% ----------------------------------------------------------

% --- COLOR SETUP ---
\usepackage{xcolor}

\definecolor{cream}{RGB}{246,242,219}
\definecolor{teamred}{RGB}{209,57,64}
\definecolor{navy}{RGB}{40,46,87}
\definecolor{softwhite}{RGB}{250,250,230}
\definecolor{accentyellow}{RGB}{240,200,60}

% Extra utility colors
\definecolor{muted}{RGB}{120,125,140}
\definecolor{darkcream}{RGB}{230,225,205}

% ----------------------------------------------------------
%  FONT SETUP (Optional but Recommended)
%  Works best with XeLaTeX or LuaLaTeX
% ----------------------------------------------------------
\usepackage{fontspec}

% Main document font (Google fonts — upload to Overleaf)
% If you don't upload fonts, comment these two lines.
\setmainfont{SourceSans3}[
  Path=./,
  ItalicFont  = *-Italic,
  BoldFont    = *-Bold
]

\newfontfamily\headingfont{ScienceGothic}[
  Path=./,
]

% Font Awesome for icons
\usepackage{fontawesome5}

% ----------------------------------------------------------
%  SECTION STYLE
% ----------------------------------------------------------
\usepackage{titlesec}

% Section
\titleformat{\section}
  {\Large\headingfont\bfseries\color{navy}}
  {\thesection}
  {1em}
  {\rule{\linewidth}{1pt}\vspace{0.5em}\\}

% Subsection
\titleformat{\subsection}
  {\large\headingfont\bfseries\color{teamred}}
  {\thesubsection}
  {0.8em}
  {}

% Subsubsection
\titleformat{\subsubsection}
  {\bfseries\color{navy}}
  {\thesubsubsection}
  {0.6em}
  {}

% ----------------------------------------------------------
%  PAGE STYLE
% ----------------------------------------------------------
\usepackage{fancyhdr}
\pagestyle{fancy}

\fancyhf{} % clear all

\lhead{\textcolor{navy}{\headingfont \courseTitle}}
\rhead{\textcolor{teamred}{\thepage}}
\renewcommand{\headrule}{\color{teamred}\hrule height 1pt}

% ----------------------------------------------------------
%  HYPERLINK STYLE
% ----------------------------------------------------------
\usepackage{hyperref}

\hypersetup{
  colorlinks = true,
  linkcolor  = teamred,
  urlcolor   = navy,
  citecolor  = accentyellow
}

% ----------------------------------------------------------
%  CUSTOM BOX ENVIRONMENTS
% ----------------------------------------------------------
\usepackage{tcolorbox}
\tcbuselibrary{skins,breakable}

% Info box
\newtcolorbox{infobox}{
  enhanced,
  breakable,
  colback   = softwhite,
  colframe  = navy,
  coltext   = navy,
  boxrule   = 1pt,
  arc       = 4pt,
  left=6pt, right=6pt, top=6pt, bottom=6pt
}

% Warning box
\newtcolorbox{warningbox}{
  enhanced,
  breakable,
  colback   = accentyellow!30,
  colframe  = accentyellow!80!navy,
  coltext   = navy,
  boxrule   = 1pt,
  arc       = 4pt,
  left=6pt, right=6pt, top=6pt, bottom=6pt
}

% Highlight box
\newtcolorbox{highlightbox}{
  enhanced,
  breakable,
  colback   = teamred!10,
  colframe  = teamred,
  coltext   = navy,
  boxrule   = 1pt,
  arc       = 4pt,
  left=6pt, right=6pt, top=6pt, bottom=6pt
}

% ----------------------------------------------------------
%  ITEMIZE / ENUMERATE STYLE
% ----------------------------------------------------------
\usepackage{enumitem}

\setlist[itemize]{
  label=\textcolor{teamred}{\large\textbullet},
  leftmargin=1.2em,
  itemsep=4pt
}

\setlist[enumerate]{
  label=\textcolor{navy}{\arabic*.},
  leftmargin=1.4em,
  itemsep=4pt
}

% ----------------------------------------------------------
%  TITLE STYLE
% ----------------------------------------------------------
\usepackage{titling}

\pretitle{
  \begin{center}
  \headingfont\bfseries\LARGE\color{navy}
}
\posttitle{
  \end{center}
}

\preauthor{
  \begin{center}
  \color{teamred}
}
\postauthor{
  \end{center}
}

% ----------------------------------------------------------
%  OPTIONAL — SOFT PAGE BACKGROUND
%  Comment out if you want white pages
% ----------------------------------------------------------
\usepackage{pagecolor}
\pagecolor{cream}
\color{navy}

% ----------------------------------------------------------
% CUSTOM COMMANDS FOR TEMPLATE
% ----------------------------------------------------------
\newcommand{\courseTitle}{Cool Stuff on Linux — A Reference Guide for Adventurous Newcomers}
\newcommand{\courseDate}{27-12-2025}
\newcommand{\authorName}{Khaled Mahfouz}
\newcommand{\contactInfo}{\texttt{@khaledmhfz}}

% Minimal icon commands
\newcommand{\bookicon}{\faBook}
\newcommand{\videoicon}{\faVideo}
\newcommand{\searchicon}{\faSearch}
\newcommand{\syncicon}{\faSync}
\newcommand{\wineicon}{\faWindows}
\newcommand{\packageicon}{\faBox}
\newcommand{\shellicon}{\faTerminal}
\newcommand{\toolicon}{\faTools}
\newcommand{\themeicon}{\faPalette}
\newcommand{\rocketicon}{\faRocket}
\newcommand{\aiicon}{\faRobot}

% ----------------------------------------------------------
% END OF THEME
% ----------------------------------------------------------

\begin{document}

% Title with minimal space
\begin{center}
\headingfont\Huge\bfseries\color{navy}
\bookicon\ \courseTitle\\[0.5em]
\Large\color{teamred}
\courseDate
\end{center}

\vspace{1em}
\begin{center}
\color{teamred}\rule{0.8\textwidth}{1pt}
\end{center}

\section*{\searchicon\ 1) Research Skills \& Problem Solving}

\textbf{Goal:} Learn how to find answers fast when the terminal throws tantrums.
\textbf{What to do:}

\begin{itemize}
    \item Search the \textbf{man pages} (\texttt{man rsync}, \texttt{man bash}) — the classic built-in Linux docs.
    \item Use \texttt{apropos}, \texttt{man -k}, \texttt{--help} flags.
    \item Stack Exchange (Unix \& Linux), distro forums and reddit are gold mines.
\end{itemize}

\begin{infobox}
\faLightbulb\ \textbf{Tips:} always check dates \& distro versions — old scripts can break what you have achieved.\\
\faArrowRight\ Unix \& Linux Q\&A sorted by votes: \textit{awesome community answers} (\href{https://learnbyexample.github.io/scripting_course/Linux_curated_resources.html}{\texttt{Learn By Example}})
\end{infobox}

\section*{\videoicon\ 2) Learn \textbf{FFmpeg}}

FFmpeg is your \textbf{multimedia Swiss Army knife} — convert, cut, encode audio \& video from the terminal.

\begin{infobox}
\bookicon\ \textbf{Tutorial:} \textit{Install and Use ffmpeg in Linux} — super clear with examples.\\
\faArrowRight\ \href{https://itsfoss.com/ffmpeg/}{\texttt{https://itsfoss.com/ffmpeg/}} (\texttt{It's FOSS})\\
\bookicon\ \textbf{Main Documentation}\\
\faArrowRight\ \href{https://ffmpeg.org/ffmpeg.html}{\texttt{https://ffmpeg.org/ffmpeg.html}}
\end{infobox}

\textbf{Basics:}

\begin{highlightbox}
\texttt{\# convert an MKV to MP4}\\
\texttt{\$ ffmpeg -i input.mkv output.mp4}
\end{highlightbox}

Fun hack: use FFmpeg to record your screen or trim clips like a wizard. \faMagic

\section*{\syncicon\ 3) FOSS Alternatives for Common Software}

Why pay or use closed source when Linux offers freedom (and bragging rights)?

\textbf{Check out:}

\begin{itemize}
    \item LibreOffice → Office suite
    \item GIMP → Photoshop replacement
    \item Inkscape → vector graphics
    \item Blender → 3D beast
    \item Audacity → audio editing
    \item Krita → painting/art
\end{itemize}

\begin{infobox}
\faArrowRight\ Neowin guide on top FOSS apps for Linux → \href{https://www.neowin.net/guides/top-10-foss-apps-to-make-your-linux-experience-more-enjoyable/}{\texttt{https://www.neowin.net/guides/...}} (\texttt{Neowin})\\
Also use \href{https://alternativeto.net}{\texttt{AlternativeTo.net}} for more suggestions.
\end{infobox}

\section*{\wineicon\ 4) Run Windows Apps/Games with Bottles, Wine, Proton \& Lutris}

Welcome to \textbf{Windows software on Linux} land!

\textbf{Tech stack:}

\begin{itemize}
    \item \textbf{Wine} — compatibility layer to run many Windows apps and games. (\href{https://en.wikipedia.org/wiki/Wine_\%28software\%29}{\texttt{Wikipedia}})
    \item \textbf{Bottles} — GUI wineprefix \& dependency manager. (\href{https://www.gamingonlinux.com/2021/12/use-wine-for-gaming-on-linux-try-out-bottles/page\%3D2/}{\texttt{GamingOnLinux}})
    \item \textbf{Proton} — Wine optimized for games, built into Steam. (\href{https://en.wikipedia.org/wiki/Proton_\%28software\%29}{\texttt{Wikipedia}})
    \item \textbf{Lutris} — game launcher \& scripts installer. (\href{https://en.wikipedia.org/wiki/Lutris}{\texttt{Wikipedia}})
\end{itemize}

\begin{infobox}
\bookicon\ \textbf{Tutorials \& resources:}
\begin{itemize}
    \item Bottles + Wine on YouTube: \href{https://youtu.be/xGbQZt96q54?si=FgQJ__YMpISs7L1c}{\texttt{https://youtu.be/xGbQZt96q54?si=FgQJ\_\_YMpISs7L1c}} (\texttt{YouTube})
    \item Lutris: \href{https://youtu.be/JXHvGjiJJG0?si=1fI81xIr7zichNDy}{\texttt{https://youtu.be/JXHvGjiJJG0?si=1fI81xIr7zichNDy}}
    \item ProtonDB (community compatibility data — must bookmark).
\end{itemize}
\end{infobox}

\begin{infobox}
\faLightbulb\ \textbf{Tip:} Steam + Proton is often easiest for games; Lutris handles other launchers \& installers.
\end{infobox}

\section*{\syncicon\ 5) Learn \textbf{rsync}}

\textbf{rsync} is a \textit{super efficient} file sync tool — perfect for backups and transfers.

\begin{infobox}
\bookicon\ \textbf{Guide:} DigitalOcean tutorial — concise, example-rich.\\
\faArrowRight\ \href{https://www.digitalocean.com/community/tutorials/how-to-use-rsync-to-sync-local-and-remote-directories}{\texttt{(DigitalOcean)}}\\
\videoicon\ \textbf{Youtube Tutorial}\\
\faArrowRight\ \href{https://youtu.be/GqSxR93xK6E?si=c2TvpU8R-53SXDty}{\texttt{Backing up a Linux Server with rsync}}
\end{infobox}

\textbf{Key command patterns:}

\begin{highlightbox}
\texttt{\# basic local sync}\\
\texttt{\$ rsync -av /src/ /dest/}\\
\\
\texttt{\# remote sync}\\
\texttt{\$ rsync -avz /src/ user@host:/dest/}
\end{highlightbox}

Learn the difference trailing slash makes — it's the "slash of power" \faCrosshairs.

\section*{\aiicon\ 6) Using AI on Linux}

AI tools are popping up everywhere — including Linux!

\begin{itemize}
    \item GPT models locally: GPT4All, Ollama, etc.
    \item Online AI models.
    \item AI helpers can summarize logs, optimize scripts, and suggest commands.
\end{itemize}

\textbf{Notable tools:}

\begin{itemize}
    \item GPT4All (local transformers)
    \item Ollama (local LLM host)
    \item Bard/ChatGPT in the browser
\end{itemize}

\section*{\toolicon\ 7) Install Common GUI Apps}

Use \texttt{apt}, \textbf{Flatpak}, or \textbf{Snap}:

\textbf{APT (Debian/Ubuntu):}

\begin{highlightbox}
\texttt{\$ sudo apt update}\\
\texttt{\$ sudo apt install vlc gnome-tweaks}
\end{highlightbox}

\textbf{Flatpak:}

\begin{highlightbox}
\texttt{\$ sudo apt install flatpak}\\
\texttt{\$ flatpak install flathub com.obsproject.Studio}
\end{highlightbox}

\textbf{Snap:}

\begin{highlightbox}
\texttt{\$ sudo apt install snapd}\\
\texttt{\$ snap install spotify}
\end{highlightbox}

\begin{infobox}
\faInfoCircle\ Snaps auto-update, are sandboxed, and sometimes bigger in size. (\href{https://linuxcommando.blogspot.com/2018/}{\texttt{Linux Commando}})
\end{infobox}

\section*{\packageicon\ 8) Flatpak \& Snap (Universal Packaging)}

\textbf{Flatpak} — universal, distro-agnostic, sandboxed.
\textbf{Snap} — Canonical's system with auto-updates and confinement.
Tutorial: Linux learning guide covers both package systems.

\textbf{Example commands:}

\begin{highlightbox}
\texttt{\$ flatpak search editor}\\
\texttt{\$ snap find browser}
\end{highlightbox}

\begin{infobox}
\faLightbulb\ \textbf{Pro tip:} \textit{Flatpak often has newer versions than distro repos.}
\end{infobox}

\section*{\shellicon\ 9) Shell Aliases \& Customizations}

Aliases save time — like macros for your shell.

\begin{highlightbox}
\texttt{\$ alias ll='ls -la'}\\
\texttt{\$ alias gs='git status'}
\end{highlightbox}

Add to \texttt{\textasciitilde/.bashrc} or \texttt{\textasciitilde/.zshrc} and \texttt{source \textasciitilde/.bashrc} to reload.

You'll thank yourself later when typing gets lazy \faSmile

\section*{\toolicon\ 10) Install \& Configure Common Software}

Here's a quick starter pack for Ubuntu:

\begin{highlightbox}
\texttt{\$ sudo apt update}\\
\texttt{\$ sudo apt install \textbackslash}\\
\texttt{    git htop tmux neofetch \textbackslash}\\
\texttt{    gufw gnome-tweaks}
\end{highlightbox}

Want a script? Store it in \texttt{setup.sh} and run \texttt{bash setup.sh}.

Also check \textbf{Perfect Ubuntu Guide} on GitHub for curated apps \& setup tips.
\begin{infobox}
\faArrowRight\ \href{https://github.com/mikeroyal/Perfect-Ubuntu-Guide}{\texttt{https://github.com/mikeroyal/Perfect-Ubuntu-Guide}} (\texttt{GitHub})
\end{infobox}

\section*{\themeicon\ 11) Fonts \& Themes}

Make your desktop \textit{pop}:

\begin{itemize}
    \item Install custom fonts (\texttt{sudo apt install fonts-roboto})
    \item GNOME Tweaks for themes
    \item Flatpak themes from Flathub
\end{itemize}

Fonts + themes = aesthetic \textasciitilde100\% more hipster points \faSmile

\begin{infobox}
\videoicon\ \textbf{Youtube Tutorial}\\
\faArrowRight\ \href{https://mmbesar.github.io/Tutorials/Best-Arabic-Font-in-Linux/}{\texttt{https://mmbesar.github.io/Tutorials/Best-Arabic-Font-in-Linux/}}
\end{infobox}

\section*{\rocketicon\ 12) Cool Command Line Utilities}

Let's finish with \textbf{fun CLI tools} you'll use all the time:

\begin{itemize}
    \item \texttt{fzf} — fuzzy finder wizard
    \item \texttt{broot} — better directory tree explorer
    \item \texttt{bat} — cat with syntax highlighting
    \item \texttt{htop} / \texttt{btop} — fancy process views
    \item \texttt{ddgr} — DuckDuckGo in terminal
\end{itemize}

See a huge curated list of utilities (bookmarks worthy).
\begin{infobox}
\faArrowRight\ \href{https://github.com/Tinram/Linux-Utilities}{\texttt{https://github.com/Tinram/Linux-Utilities}} (\texttt{GitHub})\\
\videoicon\ \textbf{Youtube Video}\\
\faArrowRight\ \href{https://youtu.be/nCS4BtJ34-o?si=T6gagbT5k71H5BNE}{\texttt{12 GREAT command line programs YOU recommended!}}
\end{infobox}

\vspace{1em}
\begin{center}
\color{teamred}\rule{0.8\textwidth}{1pt}
\end{center}

\section*{\faLinux\ Wrap-Up}

You now have a \textbf{jumpstart roadmap} for fun, practical, and beginner-friendly Linux exploration — with links to \textbf{quality tutorials} and docs. Dive in, google everything (searching is half the skill \faSmile), and don't be afraid to break stuff — \textit{that's how mastery happens}.

If you want, I can turn this into a \textbf{multi-page PDF or structured course outline} next \rocketicon.

Happy hacking! \faBrain\faLaptopCode

\section*{Helpful Links}

\begin{itemize}
    \item \href{https://learnbyexample.github.io/scripting_course/Linux_curated_resources.html}{\texttt{Linux CLI and Shell scripting | scripting\_course}}
    \item \href{https://itsfoss.com/ffmpeg/}{\texttt{Install and Use ffmpeg in Linux [Complete Guide]}}
    \item \href{https://www.neowin.net/guides/top-10-foss-apps-to-make-your-linux-experience-more-enjoyable/}{\texttt{Top 10 FOSS apps to make your Linux experience more enjoyable - Neowin}}
    \item \href{https://en.wikipedia.org/wiki/Wine_\%28software\%29}{\texttt{Wine (software)}}
    \item \href{https://www.gamingonlinux.com/2021/12/use-wine-for-gaming-on-linux-try-out-bottles/page\%3D2/}{\texttt{Use Wine for gaming on Linux? Try out Bottles | GamingOnLinux}}
    \item \href{https://en.wikipedia.org/wiki/Proton_\%28software\%29}{\texttt{Proton (software)}}
    \item \href{https://en.wikipedia.org/wiki/Lutris}{\texttt{Lutris}}
    \item \href{https://www.youtube.com/watch?v=xrWSZXJbR2I}{\texttt{How to setup Lutris \& Bottles in 2025! (UPDATED GUIDE) - YouTube}}
    \item \href{https://www.digitalocean.com/community/tutorials/how-to-use-rsync-to-sync-local-and-remote-directories}{\texttt{How To Use Rsync to Sync Local and Remote Directories}}
    \item \href{https://linuxcommando.blogspot.com/2018/}{\texttt{Linux Commando: 2018}}
    \item \href{https://github.com/mikeroyal/Perfect-Ubuntu-Guide}{\texttt{GitHub - mikeroyal/Perfect-Ubuntu-Guide}}
    \item \href{https://github.com/Tinram/Linux-Utilities}{\texttt{GitHub - Tinram/Linux-Utilities}}
\end{itemize}

% Updated Footer
\vspace{\fill}
\begin{center}
\color{muted}
All rights reserved for \texttt{TogetherTeam} [NOT For Commercial Use] !!\\
contact the author via email \href{mailto:khaledmhfz2004@gmail.com}{\texttt{khaledmhfz2004@gmail.com}}
\end{center}

\end{document}